\section{几何 Langlands 等价}\srcpage{11}

本节解释几何 Langlands 等价的大致形态。

作为一般设定，我们稍微改变记号，固定代数闭域 $k$ 上的一条光滑、紧合、连通曲线 $X$。为方便起见，固定 $X$ 上的一个 Theta 特征 $\Theta$，即 $\Theta^{\otimes 2} = \OmX$；这将是归一化几何 Langlands 等价所需要的。此外，固定 $k$ 上的一个约化群 $G$。由于 $k$ 代数闭，$G$ 总是分裂的；我们固定一个钉扎（pinning），特别地固定一个 Borel 子群与极大环面 $T \subset B \subset G$。

还有第三个选择——$k$ 上（有限型）概形的上同调理论，系数取在某个特征为 $0$ 的域 $e$ 中。\footnote{应当存在允许更一般系数的变体，但如同我们的主要参考文献，我们将只关注这一情形。}在现代的理解中，这由一个六函子形式化
\[
D \colon \Corr(\Sch_k^{\mathrm{ft}}) \to \PrL_e
\]
给出，即从对应的 $((2,1)$-)范畴到 $e$-线性可表现 $(\infty,1)$-范畴的松弛对称幺半函子；详见 Heyer–Mann（2024），其先声见 Gaitsgory–Rozenblyum（2017）、Liu–Zheng（2024）与 Mann（2022，附录 A.5）。

\begin{remark}
这里有两个不同的域 $k$ 与 $e$，我们将在二者之上做代数几何。最终我们会在 $D$-模的框架中工作，此时 $k = e$；但对函数域情形的应用，有 $k = \overline{\F}_q$，而 $e = \Ql$。
\end{remark}

\srcpage{12}目前还不知道如何对任意（合理的）上同调理论的选择表述几何 Langlands 等价。相反，对标准类型的上同调理论，例如 de Rham、Betti 与 étale，人们可以写出表述，而这些表述实际上各自用到了给定情形的某些特殊性质。出于这个原因，我不再过多停留于一般性；读者应把 $D$ 想成标准上同调理论之一：de Rham、Betti 或 étale。

第一个主要角色是 $X$ 上 $G$-丛的模空间 $\Bun_G$。它是光滑 Artin 叠，其 $S$-值点是 $X \times_k S$ 上 $G$-丛的群胚，其中 $S$ 为 $k$ 上任意概形。用下降把六函子形式化 $D$ 从概形延拓到 Artin 叠，便得到 $\Bun_G$ 上层的 $(\infty,1)$-范畴 $D(\Bun_G)$。

几何 Langlands 纲领的出发点是这样一个想法：这个范畴是自守形式向量空间的一个化身。这一想法源于 Drinfeld（1983）关于算术 Langlands 纲领函数域情形的工作。也就是说，对于 $\F_q$ 上的概形 $S$，构造 $S(\F_q)$ 上有趣函数的一般配方如下。从 $S_{\overline{\F}_q}$ 上一个带 Frobenius 结构的 $\ell$-进层 $\mathcal{F}$ 出发，即一个映射 $\phi \colon \mathrm{Fr}^*\mathcal{F} \to \mathcal{F}$。然后对每个 $\F_q$-点 $s \in S$，联系上 $\mathcal{F}_s$ 上 Frobenius 的迹
\[
\tr(\phi_s|\mathcal{F}_s) \in \Ql,
\]
其中 $\overline{s}$ 是 $s$ 上方某个几何点。事实上，最近人们注意到，Frobenius 迹这一概念在范畴层面也有意义，于是可以取
\[
\tr(\mathrm{Fr}|D(\Bun_G)) \in D(\Ql),
\]
得到一个（复形的）$\Ql$-向量空间。对层取 Frobenius 迹的通常操作给出一个映射
\[
\tr(\mathrm{Fr}|D(\Bun_G)) \to \cA_c(\Bun_G(\F_q)).
\]

下面的定理由 Arinkin、Gaitsgory、Kazhdan、Raskin、Rozenblyum 与 Varshavsky（2021）证明，它给出了「$D(\Bun_G)$ 是自守形式空间的范畴化」这一想法的一个非常强的体现：

\begin{theorem}
给 $D(\Bun_G)$ 取一个合适的定义（即取 étale $\Ql$-层构成的范畴），自然映射
\[
\tr(\mathrm{Fr}|D(\Bun_G)) \to \cA_c(\Bun_G(\F_q))
\]
是同构。
\end{theorem}

这个定理的一个关键部分是对 $D(\Bun_G)$ 的正确定义。遵循 Laumon（1987，第 6 节）的想法，必须取具有幂零奇异支撑的层，这里的幂零奇异支撑是就层的微局部理论而言的。我们将在下文第 3 节讨论这一点。

\srcpage{13}
\begin{remark}
这大概是我们可以看到的第一个场合，其中人们在几何 Langlands 纲领中做「用范畴做代数」。换句话说，大约一个世纪以前，人们逐步从用数做代数（加、乘……）走向用模做代数（直和、张量积……）；而在这里，人们做的是用范畴做代数（范畴同样容许直和、张量积……）。这样做的可能性依赖于 Lurie 的基础性工作，特别是 Lurie（2017），并要求全程以 $(\infty,1)$-范畴的方式工作。

几何 Langlands 猜想的证明展现了在用范畴做代数方面的极高技巧——所有关键步骤都在于对范畴的巧妙定义与操作（当然其间穿插着几何洞见）。值得注意的是，并且有点反直觉的是，这些操作最终能证明非常具体的东西，例如 Arthur–Ramanujan 猜想。
\end{remark}

几何 Langlands 纲领的目标是用「谱的语言」描述 $D(\Bun_G)$。粗略地说，人们遵循算术情形中的纲要：即第一步是构造一个作用在 $D(\Bun_G)$ 上的大交换算子代数。\footnote{延续上一条注，我们在这里使用「范畴中作用于一个范畴的交换代数」这一概念。注意，「范畴中的交换代数」与对称幺半范畴是同一回事。作为要记住的基本例子，对任意概形或叠 $X$，其（导出 $(\infty,1)$-）拟凝聚层范畴 $D_{\qc}(X)$ 是一个「范畴中的交换代数」；而对 $X$ 上的任意叠 $Y$，$D_{\qc}(Y)$ 成为 $D_{\qc}(X)$ 上的一个模。在 Lurie 的基础中，这并不像看起来那么可怕：首先，对任意对称幺半 $(\infty,1)$-范畴 $\mathcal{C}$，有（交换）代数 $A \in \mathcal{C}$ 的概念，以及 $A$-模 $M \in \mathcal{C}$ 的概念。这些基本上是显然的概念——$A$ 带有一个单位 $1 \to A$、一个乘法 $A \otimes A \to A$，并满足结合律与交换律。类似地，$A$-模 $M$ 带有一个作用映射 $A \otimes M \to M$，使通常的图表交换。重要的微妙之处在于，要求结合律、交换律等的图表不是条件，而是额外的数据，它们服从更高阶的相容性。现在 $\PrL_e$ 本身就是一个对称幺半 $(\infty,1)$-范畴，因此可以把上述概念也应用于 $\PrL_e$ 中的交换代数，得到其上的模等概念。}

事实上，人们可以用完全类似的方式定义 Hecke 算子。也就是说，对每个点 $x \in X$，有 Hecke 对应
\[
\begin{tikzcd}
& \Hck_{G,x} \arrow[dl, "p_1"'] \arrow[dr, "p_2"] & \\
\Bun_G & & \Bun_G,
\end{tikzcd}
\]
其中 $\Hck_{G,x}$ 参数化两个 $G$-丛以及它们在 $x$ 之外的同构。在算术情形中，Hecke 算子由元素 $h \in \mathbb{C}[G(\cO_x)\backslash G(F_x)/G(\cO_x)]$ 列举，并由如下操作给出。取 $\Bun_G(\F_q)$ 上的一个函数，把它拉回到 $\Hck_{G,x}(\F_q)$，乘以 $h$，然后沿 $p_2$ 的纤维求和。所有这些操作在层的六种运算下都有类比：拉回是 $p_1^*$，乘法\srcpage{14}是 $\otimes$，沿纤维求和是 $p_{2!}$。\footnote{事实上，在取 Frobenius 迹的层-函数对应下，这一类比是精确的。}

然而，还有一个问题有待回答：$h$ 的类比物是什么？这里需要的不是函数，而是一个层，且它应当是空间 $G(\cO_x)\backslash G(F_x)/G(\cO_x)$ 上的层，其中 $\cO_x \cong k[[t]]$。实际上，正如 $\Bun_G$，这应当被理解为一个模问题，即局部 Hecke 叠 $\Hck^{\loc}_G$，它是把 $k$-代数 $A$ 映到
\[
G(A[[t]])\backslash G(A((t)))/G(A[[t]])
\]
的函子的层化。采用记号 $\LG(A) = G(A((t)))$ 与 $\LpG(A) = G(A[[t]])$ 分别表示圈群与正圈群，这由
\[
\Hck^{\loc}_G = \LpG\backslash \LG/\LpG = \LpG\backslash \Gr_G
\]
给出，其中 $\Gr_G = \LG/\LpG$ 是称为仿射 Grassmann 簇的归纳投射概形。

正如通常的 Hecke 代数实际上是卷积下的代数，$e$-线性 $(\infty,1)$-范畴 $D(\Hck^{\loc}_G)$ 具有代数结构，即配备了一个幺半张量积。

几何 Langlands 纲领的一个关键输入是几何 Satake 等价，即 Satake 同构的范畴版本。为了陈述它，人们要过渡到反常层的子范畴。同样，这对一般的上同调理论没有明确含义，但在 de Rham、Betti 与 étale 情形中有既定含义。\footnote{然而，最近 Richarz–Scholbach（2021）的工作给出了几何 Satake 等价的一个动机版本，其办法是巧妙利用相关几何中出现的受限类型的动机——仅混合 Tate 动机。}

\begin{theorem}[Mirković–Vilonen，2007]
存在幺半范畴的自然等价
\[
\Perv(\Hck^{\loc}_G) \cong \Rep(\Gd),
\]
其中 $\Rep(\Gd)$ 是 Langlands 对偶群 $\Gd$ 的有限维表示的 Abel 范畴。
\end{theorem}

特别地，$\Perv(\Hck^{\loc}_G)$ 竟然奇迹般地是对称幺半的，这反映了通常 Hecke 代数的不显然的交换性。

\begin{remark}
人们可以在这一等价上取 Frobenius 迹，以还原通常的 Satake 同构。事实上，对于分歧群——即 $k((t))$ 上不承认 $k[[t]]$ 上约化模型的群 $G$——存在几何 Satake 等价的各种版本，可用以（重新）证明分歧群的 Hecke 代数的 Satake 同构，参见 Zhu（2024）、van den Hove（2024，定理 1.4）。\srcpage{15}这可以看作用来识别自守函数空间的那个论证的一个极小的玩具版本。
\end{remark}

我们注意到，几何 Satake 等价是几何给出比算术更精细之物的第一个关键点。也就是说，通常的 Satake 同构只把对偶群作为一个方便的记账工具引入——真正被 Satake 同构强制出来的只有 $\check T/W$。而在这里，人们可以用 Tannakian 理论，从仅涉及 $G$ 的数据真正重构出对偶群 $\Gd$。可以说，这实际上就是 Langlands 对偶群的正确界定——它就是在几何 Satake 等价中出现的 Tannakian 群。

\begin{remark}
几何 Satake 等价的证明，是先构造 $\Perv(\Hck^{\loc}_G)$ 上的交换性约束，即把幺半结构提升为对称幺半结构。这用到另一个关键想法：融合。上面我们在任意点 $x \in X$ 处引入了局部 Hecke 叠，但人们同样可以在任意有限点集 $(x_1,\ldots,x_n)$ 处定义局部 Hecke 叠。一般地，当这些点互不相同时，它就是各个点处局部 Hecke 茎的乘积。然而，当点碰撞时，相应的因子也碰撞：若所有点都是同一个点 $x$，它就只是 $x$ 处的局部 Hecke 叠。总而言之，这给出族中的一个（Beilinson–Drinfeld）局部 Hecke 叠
\[
\Hck^{\loc}_{X^n} \to X^n,
\]
它参数化点 $(x_1,\ldots,x_n) \in X^n$ 以及 $X$ 在诸 $x_i$ 之并处的完备化上的两个 $G$-丛，并带有它们在诸 $x_i$ 之外的同构。给定局部 Hecke 叠上的反常层 $\mathcal{F}_1,\ldots,\mathcal{F}_n$，人们可以在 $\Hck^{\loc}_{X^n}$ 上定义它们的融合积 $\mathcal{F}_1 * \cdots * \mathcal{F}_n$，它一般地只是诸 $\mathcal{F}_i$ 的外张量积。为把它延拓到处处，结果表明，「在 $X^n$ 上泛局部无环」这一性质是一种「平坦性」或「连续变化」假设，它保证延拓唯一，从而给出融合积。限制到对角点 $(x,\ldots,x) \in X^n$，这就还原了幺半张量积，但现在它显然是对称的。\footnote{一个技术要点是，交换性约束必须被某些符号修正，才能使几何 Satake 等价成立——这里的问题是在奇上同调次数中出现的通常的 Koszul 符号。}
\end{remark}

这一过程——让 $x \in X$ 在曲线上移动，乃至让两个点移动并碰撞——是通常算术 Langlands 纲领中看不到的新可能性。一言以蔽之，正是这种额外的灵活性使几何 Langlands 纲领更易接近。

利用几何 Satake 等价，对每个点 $x \in X$ 我们得到 $\Rep(\Gd)$ 在 $D(\Bun_G)$ 上的作用，从而得到 $\bigotimes_{x \in X}\Rep(\Gd)$ 的作用。这是\srcpage{16}无穷维 Hecke 代数
\[
\mathbb{T} = \bigotimes_p \mathbb{C}[G(\Z_p)\backslash G(\Q_p)/G(\Z_p)] \otimes Z(U(\gc))
\]
作用于自守函数向量空间的类比。

但正如我们刚刚注意到的，在这种几何情形中我们可以让点 $x \in X$ 变动。这给出族中的 Hecke 对应
\[
\begin{tikzcd}
& \Hck_{G,X} \arrow[dl] \arrow[dr] & \\
\Bun_G & & \Bun_G \times X,
\end{tikzcd}
\]
其中 $\Hck_{G,X}$ 参数化一个点 $x \in X$ 以及两个在 $x$ 之外同构的 $G$-丛。Hecke 算子现在可以定义为函子
\[
D(\Bun_G) \to D(\Bun_G \times X),
\]
对所有 $x \in X$ 插值通常的 Hecke 算子。这给出各个 $x \in X$ 的作用之间的一些非平凡相容性，并可用于产生谱作用。

为了解释谱作用，我们必须引入几何 Langlands 纲领中的另一个关键角色：$X$ 上 $\Gd$-局部系统的叠 $\LocSys_{\Gd}$，详见第 4 节。这是 $e$ 上的一个（导出）叠，参数化按所选上同调理论意义下 $X$ 上的 $\Gd$-局部系统。换言之，若所选上同调理论是 de Rham（分别地，Betti、étale），则 $\LocSys_{\Gd}$ 就是 $X$ 上 de Rham（分别地，Betti、étale）$\Gd$-局部系统的模空间。出于技术原因，必须把它视为导出代数几何中的一个对象。

有了这些，第一个定理（不是历史上第一个，而是在几何 Langlands 证明的逻辑上第一个）是谱作用的存在性。我们将在第 5 节进一步说明。

\begin{theorem}
对 $D(\Bun_G)$ 取合适的定义，存在 $D_{\qc}(\LocSys_{\Gd})$ 在 $D(\Bun_G)$ 上的一个作用，它细化 Hecke 作用。特别地，对所有 $V \in \Rep(\Gd)$ 与 $x \in X$，诱导函子
\[
D(\Bun_G) \to D(\Bun_G)
\]
由 $D_{\qc}(\LocSys_{\Gd})$ 中如下对象的作用给出：该对象对每个 $\Gd$-局部系统，联系上表示 $V$ 所给出的相伴局部系统在 $x \in X$ 处的纤维。
\end{theorem}

\begin{remark}
如果我们更仔细地写出谱作用与 Hecke 作用的相容性，这实际上会唯一确定谱作用。

这相当于说，自守形式空间从 $\mathbb{T}$-模提升为 $\LocSys_{\Gd}$ 上的模！因此，它已经给出了「自守到 Galois」的方向。
\end{remark}

\srcpage{17}
\begin{remark}
在函数域情形，Lafforgue（2018）对自守形式空间证明了这一陈述的类比：自守形式空间中出现的每个 Hecke 特征值系统，都必须来自 $X$ 上的一个 $\Gd$-局部系统。他的工作甚至允许分歧的情形。他的论证依赖于非常类似的要素——几何 Satake 等价，以及让曲线上的点移动的可能性；我们将在下文注 5.4 中描述他的结果的一个视角。V. Lafforgue 的工作把几何 Langlands 纲领与算术 Langlands 纲领重新统一起来，并直接启发了许多后续工作。

在 $D$-模情形中，谱作用的存在性需要沉重的非形式输入，这归功于 Gaitsgory（2015），他依赖于 Kac–Moody 局部化的结果，而这些结果直到最近才由 Arinkin、Beraldo、Campbell、Chen、Færgeman、Gaitsgory、Lin、Raskin 与 Rozenblyum（2024）写出。在 Betti 情形，正确的范畴 $D(\Bun_G) = D_{\Nilp}(\Bun_G)$ 由 Ben-Zvi 与 Nadler（2018）分离得到，而 Nadler 与 Yun（2019）在这一情形构造了谱作用。有点令人惊讶的是，在该设定中这本质上是形式废话；人们除了（高度结构化的）Hecke 作用之外不需要任何其他输入。他们的工作由 Arinkin、Gaitsgory、Kazhdan、Raskin、Rozenblyum 与 Varshavsky（2020b）推广到了 étale 情形。
\end{remark}

正如我们所讨论的，猜想 1.5 可以看作有两部分：首先，它预测自守形式空间提升为 $\LocSys_{\Gd}$ 上的一个模；然后它又预测这个模对应于一个具体的对象，即对偶化层。一般地，这是一个秩 $1$ 自由模。

在几何情形，定理 2.7 已给出了前半部分！现在人们希望把 $D(\Bun_G)$ 描述为 $D_{\qc}(\LocSys_{\Gd})$ 上的模。如同算术情形，最乐观的希望是 $D(\Bun_G)$ 成为 $D_{\qc}(\LocSys_{\Gd})$ 上的秩 $1$ 自由模；但这并不完全正确。相反，人们必须把由完美复形紧生成的 $D_{\qc}$ 替换为具有幂零奇异支撑的凝聚复形的归纳完备化，即 $\IndCoh_{\Nilp}(\LocSys_{\Gd})$，参见 Arinkin–Gaitsgory（2015）。

\begin{conjecture}[几何 Langlands 猜想]
存在 $D_{\qc}(\LocSys_{\Gd})$-线性 $(\infty,1)$-范畴的等价
\[
D(\Bun_G) \simeq \IndCoh_{\Nilp}(\LocSys_{\Gd}).
\]
\end{conjecture}

事实上，在这一情形中可以具体指定这个同构的归一化。也就是说，它应当把右端的结构层送到左端的 Whittaker 层。Whittaker 层的定义是 Poincaré 级数的直接类比：人们考察对应
\[
\begin{tikzcd}
& \Bun_U^{\Theta} \arrow[dl] \arrow[dr] & \\
\Bun_G & & \A^1,
\end{tikzcd}
\]
\srcpage{18}从 $\A^1$ 上的指数局部系统出发（即指数 $D$-模，分别地 Artin–Schreier 层；在 Betti 情形须作小的调整），把它拉回到 $\Bun_U^{\Theta}$，然后取到 $\Bun_G$ 的 $!$-前推。

事实证明，若存在所需形式的等价，它由这一要求唯一决定，参见下文注 9.4。再次，这相对于自守形式的情形是一个改进，见证了范畴陈述的额外刚性。

Gaitsgory–Raskin（2025b）、Arinkin、Beraldo、Campbell、Chen、Færgeman、Gaitsgory、Lin、Raskin 与 Rozenblyum（2024）、Campbell、Chen、Færgeman、Gaitsgory、Lin、Raskin 与 Rozenblyum（2024）、Arinkin、Beraldo、Chen、Færgeman、Gaitsgory、Lin、Raskin 与 Rozenblyum（2024）以及 Gaitsgory–Raskin（2026）这一系列论文，证明了特征 $0$ 的几何 Langlands 猜想。

\begin{theorem}[几何 Langlands 定理，特征 $0$]
若 $k$ 的特征为 $0$，则几何 Langlands 猜想成立，无论在 de Rham、Betti 还是 étale 情形中。
\end{theorem}

人们几乎可以消除对特征的假设：

\begin{theorem}[几何 Langlands 定理，特征 $p$，Gaitsgory–Raskin，2025a]
若 $k$ 的特征为正，且大于由 $G$ 决定的某个常数，\footnote{特别地，$p = \operatorname{char}(k)$ 不应是 $G$ 的挠素数。}则存在开闭子集 $\LocSys'_{\Gd} \subset \LocSys_{\Gd}$ 以及 $D_{\qc}(\LocSys_{\Gd})$-线性 $(\infty,1)$-范畴的等价
\[
D(\Bun_G) \simeq \IndCoh_{\Nilp}(\LocSys'_{\Gd}).
\]
\end{theorem}

正如之前所述，人们当然期望 $\LocSys'_{\Gd} = \LocSys_{\Gd}$；这对 $G = \mathrm{SL}_n$ 是已知的。

让我们以对证明策略的极简概述结束本节。首先，若 $k$ 的特征为 $0$，人们证明 de Rham、Betti 与 étale 情形中的陈述彼此等价；这是 Riemann–Hilbert 对应的一个应用，也是 Gaitsgory–Raskin（2025b）的主题。若 $k$ 的特征为正，人们通过把 $X$ 提升为 Witt 向量环 $W(k)$ 上的曲线，把 $X$ 形变到特征 $0$，并利用一般纤维上的结果，由退化推出 $X/k$ 上的结果；这是 Gaitsgory–Raskin（2025a）的主题。目前，这一论证还不能完全奏效，原因是混合特征下 étale 层的奇异支撑理论有微妙之处——这最终导致了开闭子集 $\LocSys'_{\Gd} \subset \LocSys_{\Gd}$ 的出现。

此后，证明在 de Rham 情形、$k$ 的特征为 $0$ 下进行。证明对（$G$ 的）秩作归纳，类比于把自守形式分解为 Eisenstein 级数的连续谱与尖点形式，这些 Eisenstein 级数由 $G$ 的真 Levi 子群诱导。\srcpage{19}至少有三个关键要素。

\begin{enumerate}
\item Whittaker 系数的保守性结果，归功于 Færgeman–Raskin（2025）。粗略地说，这保证 $D(\Bun_G)$ 上所有缓增的、例如尖点的层都能被检测到。
\item Kac–Moody 局部化。这给出一个非常一般的程序，从谱数据（即 opers）出发，在 $D(\Bun_G)$ 上产生（自守的）对象；它可追溯到 Beilinson–Drinfeld（1996）的开创性工作。这是仅在 $D$-模设定中存在的操作，并最终导致了「Galois 到自守」的方向。
\item $\LocSys_{\Gd}$ 的几何。对于最后的论证，即识别尖点部分，人们要对 $D$-模设定中 $\LocSys_{\Gd}$ 的几何作一些非常具体的分析。值得注意的是，与算术情形中尖点形式及其对应的 $L$-参数构成离散谱不同，在几何情形中，即便不可约 $\Gd$-局部系统的轨迹也是连通的，从而构成连续族，使插值论证得以进行。
\end{enumerate}

最后，对定理 2.12 中的等价取 Frobenius 迹，便可推出定理 1.7。这一步能还原自守形式空间，本身就是一个非平凡的陈述，参见第 8 节。

\begin{remark}
正如之前所述，几何 Langlands 猜想的一般形态，已被证明极具影响力，甚至超出了它最初的光滑紧合曲线设定。也就是说，在 Fargues–Scholze（2024）中，它在 Fargues–Fontaine 曲线上得到解释，与 $p$ 进域上的局部 Langlands 对应相联系；而在 Ben-Zvi–Nadler（2013）与 Scholze（2024）中，它在扭量 $\mathbb{P}^1$ 上得到解释（以略有不同的方式），与阿基米德域（即 $\R$ 与 $\C$）上的局部 Langlands 对应相联系。此外，$p$ 进域上局部 Langlands 对应的各种更精细形式——处理 $p$ 进向量空间上的 Banach 表示或局部解析表示——目前正用这一范式研究，例如参见 Emerton–Gee–Hellmann（2025）。这打破了该领域的一个僵局：在 $\mathrm{GL}_2(\Q_p)$ 之外，甚至猜想性的图景直到最近都还不清楚。
\end{remark}
